% ==============================================================================
% Section: Introduction
% Status: REWRITTEN (2026-05-02) — laser-focus rewrite. Lead with the puzzle:
%   binding cap, observed shortage, the consequences this paper traces.
% ==============================================================================

\section{Introduction}
\label{sec:introduction}

On January 28, 2025, Pennsylvania Governor Josh Shapiro announced an
agreement with PJM Interconnection that capped capacity prices in the
2026/27 and 2027/28 Base Residual Auctions at approximately
\$325/MW-day of unforced capacity (UCAP) and floored them at
approximately \$175/MW-day \citep{Shapiro2025_press_release,
PJM2025_settlement_slides}. The settlement followed a Section 206
complaint filed against PJM's tariff one month earlier
\citep{PJM2025_settlement_slides} and was projected by the
Commonwealth to save consumers over \$21~billion across the two
delivery years \citep{Shapiro2025_press_release}.

Standard price-control theory delivers a sharp prediction. A binding
ceiling on a market that would otherwise clear above the cap
truncates supply at the cap; if the supply willing to clear at the
cap is short of the planning target, the auction clears short.
\citet{Weitzman1974_prices} formalizes the logic in the
prices-vs-quantities framework, and
\citet{GlaeserLuttmer2003_rent} document a distinct allocative
consequence: misallocation across heterogeneous demanders,
qualitatively different from the standard deadweight triangle. The
Shapiro cap, layered on top of a publicly known administrative demand
curve and a concentrated supply, is a direct candidate for these
predictions.

Both capped auctions cleared short of the reliability requirement.
The 2026/27 BRA cleared 314~MW below the requirement (a 0.23\%
shortfall; 492~MW below the IRM-implied reserve target); the 2027/28
BRA cleared 6{,}623~MW below the requirement (a 4.7\% shortfall;
6{,}517~MW below the IRM target). The PJM Independent Market Monitor (IMM) declared both
auctions ``not competitive'' \citep{MonitoringAnalytics2025_sotm}, and
an IMM mechanical scenario that re-clears the same offers against the
unrestricted (pre-settlement) VRR curve implies that approximately
9{,}500~MW more capacity would have cleared in 2027/28---enough to
restore the reliability requirement plus a positive reserve margin.
The shortage emerged when the cap began to bind, and not before.

This paper traces the connection. Section~\ref{sec:theory} states the
textbook ceiling logic and uses the symmetric supply function
equilibrium of \citet{Klemperer1989_sfe}, \citet{Green1992_british},
and \citet{Holmberg2008_unique_sfe} to characterize the unconstrained
equilibrium price; under the IMM-reported concentration
($\mathrm{RSI}_3 < 1$), the equilibrium is at the cap and the
settlement cap binds. Section~\ref{sec:cap_incidence} documents
which administrative constraint was binding in each of the seven
completed BRAs and shows that the shortage emerged in the same years
the cap began to bind.
Section~\ref{sec:case_studies} develops the dollar and physical
consequences of the cap in each capped auction: in 2026/27 the cap
suppressed price more than quantity; in 2027/28 it suppressed both.
Section~\ref{sec:conclusion} reads the announced \$21~billion savings
claim in light of the shortfall and discusses the welfare and
investment-signal implications of a binding cap during a tight-supply
episode.
